\documentclass[a4paper,12pt]{article}

\usepackage[T2A]{fontenc}
\usepackage[utf8]{inputenc}
\usepackage[russian]{babel}
\usepackage{amsmath,amssymb}
\usepackage{PTSans}
\renewcommand{\familydefault}{\sfdefault}
\usepackage{icomma}
\usepackage[
  left=20mm,
  right=20mm,
  top=18mm,
  bottom=20mm,
  headheight=15pt
]{geometry}
\usepackage{microtype}
\usepackage{tikz}
\usetikzlibrary{calc,arrows.meta,positioning,shapes.geometric}
\usepackage{pgfplots}
\pgfplotsset{compat=1.18}
\usepackage{tabularx,booktabs}
\usepackage{enumitem}
\usepackage{parskip}
\usepackage{needspace}
\usepackage{fancyhdr}
\usepackage{setspace}
\usepackage{xcolor}
\usepackage{hyperref}

% -------------------------------------------------
% Цветовая палитра
% -------------------------------------------------

\definecolor{accent}{HTML}{008080}
\definecolor{darkaccent}{HTML}{004D4D}
\definecolor{textgray}{HTML}{333333}
\definecolor{softgray}{HTML}{777777}
\definecolor{gridgray}{HTML}{D7D7D7}

\color{textgray}
\onehalfspacing
\raggedbottom

\hypersetup{
  colorlinks=true,
  linkcolor=darkaccent,
  urlcolor=darkaccent,
  pdfauthor={Лёвин Артём Александрович},
  pdftitle={Метод узловых точек. Чек-лист для Кристины}
}

% -------------------------------------------------
% Оформление страниц
% -------------------------------------------------

\pagestyle{fancy}
\fancyhf{}
\fancyhead[L]{\small\color{softgray}Кристина}
\fancyhead[R]{\small\color{softgray}10 класс • подготовка к ЕГЭ}
\fancyfoot[C]{\small\color{softgray}\thepage}
\renewcommand{\headrulewidth}{0.3pt}
\renewcommand{\footrulewidth}{0pt}

\setlist[itemize]{
  leftmargin=1.7em,
  itemsep=0.3em,
  topsep=0.35em
}

\setlist[enumerate]{
  leftmargin=2em,
  itemsep=0.65em,
  topsep=0.5em
}

\newcommand{\topic}[1]{%
  \Needspace{8\baselineskip}%
  \par\vspace{0.8em}
  {\large\bfseries\color{darkaccent}#1}\par
  \vspace{0.2em}
  {\color{accent}\rule{\linewidth}{0.6pt}}\par
  \vspace{0.25em}
}

\newcommand{\keyidea}[1]{%
  \par\smallskip
  {\bfseries\color{darkaccent}#1}\enspace
}

% -------------------------------------------------
% Стили блок-схем
% -------------------------------------------------

\tikzset{
  flowstart/.style={
    ellipse,
    draw=darkaccent,
    line width=0.9pt,
    minimum width=3.5cm,
    minimum height=0.85cm,
    align=center,
    inner sep=4pt,
    font={\small\hyphenpenalty=10000\exhyphenpenalty=10000}
  },
  flowaction/.style={
    rectangle,
    rounded corners=3pt,
    draw=accent,
    line width=0.9pt,
    text width=7.2cm,
    minimum height=0.95cm,
    align=center,
    inner sep=5pt,
    font={\small\hyphenpenalty=10000\exhyphenpenalty=10000}
  },
  flowdecision/.style={
    diamond,
    aspect=2.7,
    draw=darkaccent,
    line width=0.9pt,
    text width=3.8cm,
    align=center,
    inner sep=1.5pt,
    font={\small\hyphenpenalty=10000\exhyphenpenalty=10000}
  },
  flowbranch/.style={
    rectangle,
    rounded corners=3pt,
    draw=accent,
    line width=0.9pt,
    text width=4.5cm,
    minimum height=1.05cm,
    align=center,
    inner sep=5pt,
    font={\small\hyphenpenalty=10000\exhyphenpenalty=10000}
  },
  flowarrow/.style={
    -{Stealth[length=3mm,width=2mm]},
    line width=0.9pt,
    draw=textgray
  }
}

\begin{document}

% =================================================
% ТИТУЛЬНЫЙ ЛИСТ
% =================================================

\hypersetup{pageanchor=false}
\begin{titlepage}
\thispagestyle{empty}
\centering

\vspace*{2.4cm}

{\Large\bfseries\color{darkaccent}Чек-лист\par}

\vspace{0.9cm}

{\huge\bfseries Метод узловых точек\par}

\vspace{0.35cm}

{\Large Восстановление формулы функции по графику\par}

\vspace{1.2cm}

{\large Линейная и квадратичная функции\par}

\vspace{1.8cm}

{\large 10 класс • подготовка к ЕГЭ\par}

\vfill

{\large\bfseries
Лёвин Артём Александрович\\
эксклюзивно для Кристины\par}

\vspace{0.7cm}

{\large Дата: \today\par}

\vspace{1.8cm}

\begin{tikzpicture}[remember picture,overlay]
  \draw[
    accent!65,
    line width=1.4pt
  ]
  ($(current page.south west)+(18mm,17mm)$)
  .. controls
  ($(current page.south west)+(58mm,31mm)$)
  and
  ($(current page.south east)+(-66mm,8mm)$)
  ..
  ($(current page.south east)+(-18mm,20mm)$);
\end{tikzpicture}

\end{titlepage}
\hypersetup{pageanchor=true}

% =================================================
% ОСНОВНОЙ ЧЕК-ЛИСТ
% =================================================

\section*{\color{darkaccent}Основной чек-лист}

\topic{1. Главная идея}

Если общий вид функции известен, то точки её графика позволяют найти неизвестные коэффициенты. Любая точка \(A(x_0;y_0)\), принадлежащая графику, удовлетворяет равенству
\[
y_0=f(x_0).
\]
Следовательно, каждая точно считанная точка даёт уравнение относительно неизвестных коэффициентов.

\keyidea{Основной принцип.}
Нужно получить столько уравнений, чтобы коэффициенты определялись однозначно. Для линейной функции \(f(x)=kx+b\) достаточно двух точек с разными абсциссами.

\topic{2. Узловые точки}

В задачах этого типа будем называть \textbf{узловой} точку графика, у которой обе координаты являются целыми числами.

Например,
\[
A(-1;-3),\qquad B(3;4)
\]
--- узловые точки, а точка \(C(2;2{,}4)\) узловой не является.

\keyidea{Почему они удобны.}
Целые координаты легче считать с рисунка точно, без приближений.

\keyidea{Как выбирать.}
Если подходящих точек несколько, при прочих равных удобнее брать точки с \(x=0\), \(y=0\), \(x=1\), \(x=-1\) или просто с небольшими по модулю координатами. Это упрощает вычисления, но не является обязательным условием.

\topic{3. Как читать координаты}

Координаты точки записываются в порядке
\[
(x;y).
\]
Первая координата --- значение \(x\), вторая --- значение \(y=f(x)\).

Например,
\[
A(-3;2)
\qquad\Longrightarrow\qquad
x=-3,\quad f(-3)=2.
\]

\topic{4. Линейная функция: восстановление формулы}

Пусть
\[
f(x)=kx+b
\]
и график проходит через точки
\[
A(-1;-3),\qquad B(3;4).
\]

Подставляем координаты обеих точек:
\[
\begin{cases}
-3=-k+b,\\
4=3k+b.
\end{cases}
\]

Из первого уравнения
\[
b=k-3.
\]
Подставляем во второе:
\[
4=3k+(k-3),
\]
поэтому
\[
4k=7,\qquad k=\frac74.
\]
Тогда
\[
b=\frac74-3=-\frac54.
\]

Итак,
\[
\boxed{f(x)=\frac74x-\frac54}.
\]

Если требуется найти \(f(-5)\), подставляем \(x=-5\):
\[
f(-5)=\frac74\cdot(-5)-\frac54=-\frac{40}{4}=-10.
\]

\keyidea{Самопроверка.}
Подставьте в найденную формулу координаты исходной точки. Например,
\[
f(3)=\frac{21}{4}-\frac54=4,
\]
что совпадает с ординатой точки \(B(3;4)\).

\topic{5. Метод подстановки при решении системы}

Рассмотрим ещё один набор точек:
\[
A(-3;-3),\qquad B(1;2).
\]
Для функции \(f(x)=kx+b\) получаем
\[
\begin{cases}
-3=-3k+b,\\
2=k+b.
\end{cases}
\]

Из первого уравнения:
\[
b=3k-3.
\]
Подставляем это выражение во второе:
\[
2=k+(3k-3),
\]
откуда
\[
4k=5,\qquad k=\frac54.
\]
Затем
\[
b=2-\frac54=\frac34.
\]
Следовательно,
\[
\boxed{f(x)=\frac54x+\frac34}.
\]

Если требуется найти \(f(-9)\), то
\[
f(-9)=\frac54\cdot(-9)+\frac34=-\frac{42}{4}=-\frac{21}{2}.
\]

\topic{6. Два разных типа вопроса}

\keyidea{Найти значение функции.}
Если дана функция
\[
f(x)=2x-1
\]
и требуется найти \(f(-5)\), число \(-5\) подставляется вместо \(x\):
\[
f(-5)=2\cdot(-5)-1=-11.
\]

\keyidea{Найти аргумент.}
Если известно, что
\[
f(x)=-5,
\]
то число \(-5\) подставляется вместо значения функции:
\[
-5=2x-1.
\]
Решаем уравнение:
\[
2x=-4,\qquad x=-2.
\]

Удобная памятка:
\[
\boxed{f(a)\;\Longrightarrow\;x=a}
\qquad\text{и}\qquad
\boxed{f(x)=c\;\Longrightarrow\;\text{решаем }f(x)=c\text{ относительно }x}.
\]

Для линейной функции с \(k\ne0\) такое уравнение имеет одно решение. Для квадратичной функции решений может быть два, одно или ни одного.

\topic{7. Быстрая проверка различия}

Пусть
\[
f(x)=3-2x.
\]
Тогда
\[
f(0)=3.
\]

Если же требуется найти \(x\), при котором \(f(x)=0\), получаем
\[
0=3-2x,
\qquad
2x=3,
\qquad
x=\frac32.
\]

Это две разные постановки задачи: в первой известно значение аргумента, во второй --- значение функции.

% =================================================
% БЛОК-СХЕМА
% =================================================

\clearpage
\thispagestyle{fancy}

\begin{center}
{\Large\bfseries\color{darkaccent}
Алгоритм: восстановление функции по графику}
\end{center}

\vspace{4mm}

\begin{center}
\begin{tikzpicture}[node distance=5.5mm and 12mm]

\node[flowstart] (start) {Начало};

\node[flowaction, below=of start] (form)
{Запишите общий вид функции и отметьте неизвестные коэффициенты};

\node[flowaction, below=of form] (points)
{Выберите хорошо читаемые точки графика. Их должно быть достаточно для определения всех неизвестных коэффициентов};

\node[flowaction, below=of points] (coords)
{Запишите координаты каждой точки в порядке \((x;y)\), где \(y=f(x)\)};

\node[flowaction, below=of coords] (system)
{Подставьте координаты в формулу и составьте систему уравнений};

\node[flowaction, below=of system] (solve)
{Решите систему и найдите коэффициенты};

\node[flowdecision, below=7mm of solve] (unique)
{Все коэффициенты найдены однозначно?};

\node[flowbranch, right=8mm of unique] (other)
{Нет: проверьте точки и уравнения; при необходимости выберите другую или дополнительную точку};

\node[flowaction, below=7mm of unique] (write)
{Да: запишите конкретную формулу функции};

\node[flowdecision, below=7mm of write] (question)
{Что требуется найти?};

\node[flowbranch, below left=9mm and 4mm of question] (value)
{\(f(a)\): подставьте \(x=a\) и вычислите значение функции};

\node[flowbranch, below right=9mm and 4mm of question] (argument)
{\(x\), если \(f(x)=c\): составьте уравнение \(f(x)=c\) и найдите все допустимые решения};

\node[flowstart, below=35mm of question] (finish) {Ответ};

\draw[flowarrow] (start) -- (form);
\draw[flowarrow] (form) -- (points);
\draw[flowarrow] (points) -- (coords);
\draw[flowarrow] (coords) -- (system);
\draw[flowarrow] (system) -- (solve);
\draw[flowarrow] (solve) -- (unique);

\draw[flowarrow]
(unique.east) -- node[above,font=\small]{нет} (other.west);

\draw[flowarrow]
(other.north) |- (points.east);

\draw[flowarrow]
(unique.south) -- node[right,font=\small]{да} (write.north);

\draw[flowarrow] (write) -- (question);

\draw[flowarrow]
(question.south west) -| node[pos=0.30,left,font=\small]{\(f(a)\)} (value.north);

\draw[flowarrow]
(question.south east) -| node[pos=0.30,right,font=\small]{\(f(x)=c\)} (argument.north);

\draw[flowarrow] (value.south) |- (finish.west);
\draw[flowarrow] (argument.south) |- (finish.east);

\end{tikzpicture}
\end{center}

\clearpage

% =================================================
% КВАДРАТИЧНАЯ ФУНКЦИЯ
% =================================================

\section*{\color{darkaccent}Продолжение чек-листа}

\topic{8. Тот же метод для квадратичной функции}

Пусть
\[
f(x)=2x^2+bx+c.
\]
Коэффициент при \(x^2\) известен, поэтому требуется найти только \(b\) и \(c\).

Пусть на графике видны точки
\[
C(0;-4),\qquad D(1;1).
\]

Для точки \(C\):
\[
-4=2\cdot0^2+b\cdot0+c,
\]
поэтому
\[
c=-4.
\]

Для точки \(D\):
\[
1=2\cdot1^2+b\cdot1+c.
\]
Подставляем \(c=-4\):
\[
1=2+b-4,
\qquad
b=3.
\]

Следовательно,
\[
\boxed{f(x)=2x^2+3x-4}.
\]
Например,
\[
f(-5)=2\cdot25+3\cdot(-5)-4=31.
\]

\keyidea{Почему точка с \(x=0\) особенно удобна.}
Для функции \(f(x)=ax^2+bx+c\)
\[
f(0)=c.
\]
Поэтому такая точка сразу даёт свободный коэффициент.

\topic{9. Сколько точек требуется}

Ориентируйтесь на число неизвестных коэффициентов и на то, позволяют ли полученные уравнения определить их однозначно.

\begin{itemize}
\item Для \(f(x)=kx+b\) неизвестны \(k\) и \(b\): нужны две точки с разными значениями \(x\).
\item Для \(f(x)=2x^2+bx+c\) неизвестны \(b\) и \(c\): достаточно двух точек с разными значениями \(x\).
\item Для \(f(x)=ax^2+bx+c\), если неизвестны все три коэффициента, нужны три точки с попарно различными значениями \(x\).
\end{itemize}

Для квадратичного многочлена три точки с различными абсциссами действительно определяют коэффициенты \(a\), \(b\), \(c\) однозначно.

\topic{10. Последний шаг после нахождения коэффициентов}

Сначала всегда запишите найденную функцию целиком. Например, если
\[
k=\frac38,
\qquad
b=-\frac12,
\]
то
\[
f(x)=\frac38x-\frac12.
\]

Затем внимательно прочитайте вопрос:
\[
f(-7),
\qquad
f(x)=5,
\qquad
f(x)=0.
\]
В первом случае подставляется значение \(x\), во втором и третьем решается уравнение относительно \(x\).

\topic{11. Типичные ошибки}

\begin{itemize}
\item Путают порядок координат: верная запись --- \((x;y)\).
\item Берут точку, которая не лежит на графике или координаты которой нельзя считать точно.
\item Путают задачи \(f(a)=?\) и \(f(x)=a\).
\item Теряют знак минус при переносе слагаемого или при подстановке отрицательного \(x\).
\item Находят один коэффициент и забывают определить остальные.
\item После решения системы не записывают итоговую формулу функции.
\item Не проверяют найденную формулу по одной из исходных точек.
\end{itemize}

Например,
\[
-3k+b=-3
\qquad\Longrightarrow\qquad
b=3k-3.
\]
При переносе \(-3k\) в правую часть его знак меняется.

\topic{12. Область определения и ограничения}

Линейная и квадратичная функции являются многочленами, поэтому определены при всех действительных \(x\):
\[
D(f)=\mathbb{R}.
\]

В рассмотренных формулах нет знаменателей с переменной, корней чётной степени или логарифмов, поэтому дополнительных ограничений на \(x\) не возникает.

% =================================================
% ТРЕНИРОВОЧНЫЙ БЛОК
% =================================================

\clearpage
\section*{\color{darkaccent}Тренировочный блок • Путь А}

\textbf{Цель:} закрепить чтение координат, восстановление формулы функции и различение задач на поиск \(f(a)\) и на поиск \(x\) по известному значению функции.

\begin{enumerate}

\item
Какие из следующих точек являются узловыми?
\[
A(-2;3),\qquad
B(1{,}5;4),\qquad
C(0;-5),\qquad
D\left(7;\frac12\right),\qquad
E(3;4).
\]

\item
Дана функция
\[
f(x)=4x+1.
\]
Найдите:
\[
\text{а) }f(1);
\qquad
\text{б) }x,\ \text{если }f(x)=1.
\]

\item
На рисунке изображён график линейной функции
\[
f(x)=kx+b.
\]
Восстановите формулу функции и найдите \(f(-5)\).

\begin{center}
\begin{tikzpicture}
\begin{axis}[
  width=0.72\linewidth,
  height=6.8cm,
  axis lines=middle,
  xmin=-4.5,
  xmax=3.5,
  ymin=-6,
  ymax=6,
  xtick={-4,-3,...,3},
  ytick={-6,-5,...,6},
  grid=both,
  minor tick num=0,
  grid style={line width=0.25pt,draw=gridgray},
  axis line style={line width=0.8pt},
  tick label style={font=\small},
  xlabel={$x$},
  ylabel={$y$},
  unit vector ratio=1 1,
  samples=200,
  clip=false
]
\addplot[
  domain=-4.3:3.3,
  line width=1.4pt,
  color=darkaccent
]
{1.5*x+0.5};

\addplot[
  only marks,
  mark=*,
  mark size=2.1pt,
  color=accent
]
coordinates {(-3,-4) (1,2)};
\end{axis}
\end{tikzpicture}
\end{center}

\item
График линейной функции
\[
f(x)=kx+b
\]
проходит через точки
\[
A(-3;-3),\qquad B(1;2).
\]
Найдите формулу функции и значение \(f(-9)\).

\newpage
\item
График линейной функции
\[
f(x)=kx+b
\]
проходит через точки
\[
A(-4;-2),\qquad B(4;1).
\]
Восстановите формулу функции и найдите \(x\), если
\[
f(x)=\frac52.
\]

\item
На рисунке изображён график функции
\[
f(x)=2x^2+bx+c.
\]
Найдите \(b\), \(c\) и вычислите \(f(2)\).

\begin{center}
\begin{tikzpicture}
\begin{axis}[
  width=0.72\linewidth,
  height=7.3cm,
  axis lines=middle,
  xmin=-3.2,
  xmax=2.3,
  ymin=-6,
  ymax=8,
  xtick={-3,-2,...,2},
  ytick={-6,-5,...,8},
  grid=both,
  minor tick num=0,
  grid style={line width=0.25pt,draw=gridgray},
  axis line style={line width=0.8pt},
  tick label style={font=\small},
  xlabel={$x$},
  ylabel={$y$},
  unit vector ratio=1 1,
  samples=250,
  clip=false
]
\addplot[
  domain=-3:2.1,
  line width=1.4pt,
  color=darkaccent
]
{2*x^2+3*x-4};

\addplot[
  only marks,
  mark=*,
  mark size=2.1pt,
  color=accent
]
coordinates {(0,-4) (1,1) (-2,-2)};
\end{axis}
\end{tikzpicture}
\end{center}

\item
График функции
\[
f(x)=-2x^2+bx+c
\]
проходит через точки
\[
D(2;5),\qquad E(3;3).
\]
Найдите коэффициенты \(b\) и \(c\), запишите формулу функции и вычислите \(f(-1)\).

\end{enumerate}

% =================================================
% ОТВЕТЫ
% =================================================

\clearpage
\section*{\color{darkaccent}Ответы}

\renewcommand{\arraystretch}{1.45}

\begin{table}[ht]
\centering
\begin{tabularx}{\linewidth}{
  >{\centering\arraybackslash}p{1.1cm}
  X
}
\toprule
\textbf{№} & \textbf{Ответ} \\
\midrule
1 & \(A,\ C,\ E\). \\
2 & а) \(5\); \qquad б) \(x=0\). \\
3 & \(\displaystyle f(x)=\frac32x+\frac12,\qquad f(-5)=-7.\) \\
4 & \(\displaystyle f(x)=\frac54x+\frac34,\qquad f(-9)=-\frac{21}{2}.\) \\
5 & \(\displaystyle f(x)=\frac38x-\frac12,\qquad x=8.\) \\
6 & \(\displaystyle b=3,\qquad c=-4,\qquad f(2)=10.\) \\
7 & \(\displaystyle b=8,\qquad c=-3,\qquad f(x)=-2x^2+8x-3,\qquad f(-1)=-13.\) \\
\bottomrule
\end{tabularx}
\end{table}

\vfill

\begin{center}
{\small\color{softgray}
Лёвин Артём Александрович эксклюзивно для Кристины}
\end{center}

\end{document}